\subsection{Beam Splitters and the Mach-Zehnder Interferometer}\label{sec:from-wave-particle-duality-to-quantum-states:beam-splitters-and-the-mach-zehnder-interferometer}

\subsubsection{The single beam splitter}\label{sec:from-wave-particle-duality-to-quantum-states:beam-splitters-and-the-mach-zehnder-interferometer:single-beam-splitter}

Until now, we have mainly studied electrons. We observed that their individual detection positions were unpredictable, even though many repeated detections produced a reproducible diffraction pattern. Now we will \hl{investigate the same underlying problem using \textbf{photons}}, the particles associated with light. Our central question remains: \textbf{\emph{what happens when we send just one quantum particle?}}

\highspace
\textcolor{Green3}{\faIcon{question-circle} \textbf{What is a beam splitter?}} A \definition{Beam Splitter} is an optical component that \hl{divides an incoming light beam} (i.e., a stream of photons) into \hl{two outgoing beams.} It can be implemented, for example, using a partially reflective mirror.

\highspace
For an ideal \textbf{50/50 beam splitter}:
\begin{itemize}
    \item 50\% of the incident light intensity is \textbf{transmitted}.
    \item 50\% of the incident light intensity is \textbf{reflected}.
\end{itemize}
So far, nothing surprising, because ordinary classical wave optics already explains how light can be divided into two beams. But now let's change the experiment.

\highspace
\textcolor{Green3}{\faIcon{question-circle} \textbf{What if we send only one photon?}} Instead of sending an intense beam of light, we \hl{consider individual photons entering the beam splitter.} We place two detectors at the outputs:
\begin{itemize}
    \item $D_1$ measures the \textbf{transmitted output}.
    \item $D_2$ measures the \textbf{reflected output}.
\end{itemize}
Each detector can produce a click (i.e., a signal) when it detects a photon.

\begin{figure}[!htp]
    \centering
    \tikzsetnextfilename{beam-splitter-single-photon-detailed}
    \begin{tikzpicture}[>=Stealth, font=\small,
            photon/.style={->, thick, decorate,
            decoration={snake, amplitude=1.2pt, segment length=7pt, post length=3mm}},
            detector/.style={draw=black!70, thick, rounded corners=2pt,
            minimum size=0.9cm, fill=gray!15},
            note/.style={draw=black!40, rounded corners=3pt, fill=black!3,
        font=\footnotesize, align=left, inner sep=4pt}]

        % beam splitter: glass slab + partially reflective coating on the incoming face
        \begin{scope}[rotate around={45:(0,0)}]
            \filldraw[fill=blue!15, draw=blue!60!black] (-0.8,-0.08) rectangle (0.8,0.08);
            \draw[blue!70!black, line width=1.4pt] (-0.8,0.08) -- (0.8,0.08);
        \end{scope}
        \node[below, font=\footnotesize, text=blue!60!black] at (0.2,-0.7) {beam splitter};

        % coating label (bottom left, arrow ends on the coated line)
        \node[font=\footnotesize, text=blue!60!black, anchor=east, align=right] at (-1.5,-0.9)
        {partially reflective coating};
        \draw[->, thin, blue!60!black] (-1.45,-0.7) -- (-0.5,-0.4);

        % input photon
        \draw[photon, red!70!black] (-4.0,0) -- (-0.12,0);
        \node[above, red!70!black] at (-2.4,0.12) {incident single photon};

        % possible outputs (lighter)
        \draw[photon, red!70!black!55] (0.12,0) -- (3.2,0);
        \draw[photon, red!70!black!55] (0,0.12) -- (0,2.25);

        % path labels
        \node[below, font=\footnotesize, text=black!70] at (1.9,-0.05) {transmitted};
        \node[left, font=\footnotesize, text=black!70] at (-0.08,1.4) {reflected};

        % detectors
        \node[detector] (D1) at (3.7,0) {$D_1$};
        \node[detector] (D2) at (0,2.7) {$D_2$};

        % probabilities
        \node[above] at (1.9,0.1) {$P(D_1)=\tfrac12$};
        \node[right] at (0.12,1.5) {$P(D_2)=\tfrac12$};

    \end{tikzpicture}
    \caption{\textbf{A single photon at a 50/50 beam splitter.} The two outputs are equally likely detection outcomes, and each photon is detected at exactly one detector.}
    \label{fig:beam-splitter-single-photon-detailed}
\end{figure}

\begin{figure}[!htp]
    \centering
    \tikzsetnextfilename{beam-splitter-cube-and-coating}
    \begin{tikzpicture}[>=Stealth, scale=0.74,
            beam/.style={very thick, red!70!black, -{Stealth[length=3mm]}},
            lab/.style={font=\small, inner sep=2pt},
        title/.style={font=\small\itshape}]

        %%%%%%%%%%%%%%%% (a) cube, side view %%%%%%%%%%%%%%%%
        \fill[blue!8]  (0,0) -- (3,0) -- (3,3) -- cycle;
        \fill[blue!16] (0,0) -- (3,3) -- (0,3) -- cycle;
        \draw[thick, black!60] (0,0) rectangle (3,3);
        \draw[line width=2.2pt, gray!65!black] (0,0) -- (3,3);
        \node[lab, gray!40!black] at (0.75,2.25) {prism};
        \node[lab, gray!40!black] at (2.35,0.4) {prism};
        \draw[gray!50!black] (0.9,0) arc[start angle=0, end angle=45, radius=0.9];
        \node[lab, gray!40!black] at (1.25,0.3) {$45^\circ$};

        \draw[beam] (-2.9,1.5) -- (1.5,1.5);
        \draw[beam] (1.5,1.5) -- (3.9,1.5);
        \draw[beam] (1.5,1.5) -- (1.5,4.6);
        \fill (1.5,1.5) circle (1.8pt);
        \node[lab, above] at (-1.75,1.55) {incident light};
        \node[lab, anchor=west, align=left] at (3.95,1.5) {transmitted\\(50\%)};
        \node[lab, above] at (1.5,4.6) {reflected (50\%)};
        \node[title] at (1.5,-0.6) {(a) side view of the cube};

        % zoom circle + link lines to the corners of panel (b)
        \draw[dashed, gray!50!black] (1.5,1.5) circle (0.5);
        \draw[dashed, gray!60!black] (1.85,1.85) -- (6.9,4.6);
        \draw[dashed, gray!60!black] (1.85,1.15) -- (6.9,0);

        %%%%%%%%%%%%%%%% (b) close-up of the coating %%%%%%%%%%%%%%%%
        \begin{scope}[shift={(9.2,2.3)}]
            \begin{scope}
                \clip (-2.3,-2.3) rectangle (2.3,2.3);
                \fill[blue!8]  (-3,-3) -- (3,3) -- (3,-3) -- cycle;
                \fill[blue!16] (-3,-3) -- (3,3) -- (-3,3) -- cycle;
                \fill[gray!65!black] (-3,-3.17) -- (3,2.83) -- (3,3.17) -- (-3,-2.83) -- cycle;
            \end{scope}
            \draw[thick, black!60] (-2.3,-2.3) rectangle (2.3,2.3);

            \node[lab, gray!40!black] at (-1.55,1.85) {glass};
            \node[lab, gray!40!black] at (1.55,-1.85) {glass};

            \draw[beam] (-2.3,0) -- (0,0);
            \draw[beam] (0,0) -- (2.3,0);
            \draw[beam] (0,0) -- (0,2.3);
            \fill (0,0) circle (1.8pt);

            \node[lab, above] at (-1.3,0.08) {$I$};
            \node[lab, below] at (1.3,-0.08) {$I/2$};
            \node[lab, left]  at (-0.08,1.25) {$I/2$};

            \node[lab, align=center, anchor=south] at (1.3,2.6)
            {thin partially reflective\\film (nanometres)};
            \draw[->, gray!40!black] (1.6,2.62) -- (1.55,1.6);

            \node[title] at (0,-2.9) {(b) close-up of the coating (not to scale)};
        \end{scope}
    \end{tikzpicture}
    \caption{%
        \textbf{How a beam splitter works.}
        (a)~A beam-splitter cube is made of two glass prisms joined along their diagonal.
        The light hits the interface at $45^\circ$: half of the intensity is transmitted straight on,
        and half is reflected at $90^\circ$.
        (b)~Close-up of the interface: the ``partially reflective mirror'' is a very thin coating,
        only a few nanometres thick, which reflects part of the light and lets the rest through.
        For a 50/50 beam splitter, the coating is tuned so that an incident intensity $I$
    leaves as $I/2 + I/2$.}
    \label{fig:beam-splitter-cube-and-coating}
\end{figure}

\noindent
\emph{\textbf{What do we observe?}} For an ideal single-photon experiment with perfect detectors, one photon produces \textbf{one detection event}: either at $D_1$ or at $D_2$. Thus, the probabilities are:
\begin{equation*}
    P(D_1)=\frac{1}{2},\qquad P(D_2)=\frac{1}{2}
\end{equation*}
However, notice the difference between two statements:
\begin{itemize}
    \item \textbf{With a classical light beam}, 50/50 refers to how the incident intensity is distributed between the outputs.
    \item \textbf{With a single photon}, 50/50 refers to the probabilities of the two possible detection outcomes.
\end{itemize}
We do \textbf{not} detect half a photon at $D_1$ and half a photon at $D_2$. So individual photon detections occur at one detector rather than being divided into fractional detection events.

\begin{deepeningbox}[: Particle indivisibility in detection]
    \textbf{Classic light beam.} Imagine sending a continuous light beam with intensity $I$ toward a 50/50 beam splitter. The two outputs have intensities:
    \begin{equation*}
        I_1 = \frac{I}{2},\qquad I_2 = \frac{I}{2}
    \end{equation*}
    Both detectors receive lights simultaneously.

    \highspace
    \textbf{Single photon.} Now imagine sending \textbf{exactly one photon} toward the beam splitter. There are two possible detection outcomes:
    \begin{itemize}
        \item $D_1$ clicks and $D_2$ does not.
        \item $D_2$ clicks and $D_1$ does not.
    \end{itemize}
    For an ideal experiment, each outcome has a probability of 50\%. What we do \textbf{not} observe is a single photon depositing half its energy at each detector. If the photon's energy is $E=h\nu$, then an ideal detection registers the photon's entire energy:
    \begin{equation*}
        \begin{aligned}
            D_1 &: E,\qquad D_2:0\\
            \text{or}\qquad D_1 &:0,\qquad D_2:E
        \end{aligned}
    \end{equation*}
    \textbf{Not}:
    \begin{equation*}
        D_1:\frac{E}{2},
        \qquad
        D_2:\frac{E}{2}
    \end{equation*}
    This is what we mean by \textbf{particle indivisibility in detection}.

    \highspace
    But \hl{this does not mean the photon necessarily travels along one definite classical path.} We are describing what happens \emph{when it is detected}, not asserting what its trajectory was before detection.
\end{deepeningbox}

\highspace
\textcolor{Green3}{\faIcon{question-circle} \textbf{What if we repeat the experiment?}} Imagine sending $1'000$ photons, one at a time. A possible sequence of detector outcomes is:
\begin{equation*}
    D_1,\ D_2,\ D_2,\ D_1,\ D_1,\ D_2,\ldots
\end{equation*}
We \textbf{cannot} generally \hl{predict which detector will register the next photon.} But after sufficiently many repetitions, approximately half of the photons will be detected at each output:
\begin{equation*}
    \frac{N_1}{N}\approx 0.5,
    \qquad
    \frac{N_2}{N}\approx 0.5
\end{equation*}
This should sound familiar! It is exactly the distinction we established in previous sections:
\begin{equation*}
    \text{unpredictable individual detections}
    +
    \text{predictable statistics}
\end{equation*}

\newpage

\begin{flushleft}
    \textcolor{Green3}{\faIcon{book} \textbf{Anticorrelated detector clicks and particle indivisibility}}
\end{flushleft}
Let's examine the detection events more closely. Imagine sending photons into the beam splitter one at a time. We observe the following possible sequence:

\begin{center}
    \begin{tabular}{@{} c c c @{}}
        \toprule
        \textbf{Photon} & \textbf{Detector $D_1$} & \textbf{Detector $D_2$} \\
        \midrule
        1 & \textcolor{Green3}{\faIcon{check}} Click & -- \\
        2 & -- & \textcolor{Green3}{\faIcon{check}} Click \\
        3 & -- & \textcolor{Green3}{\faIcon{check}} Click \\
        4 & \textcolor{Green3}{\faIcon{check}} Click & -- \\
        5 & \textcolor{Green3}{\faIcon{check}} Click & -- \\
        \bottomrule
    \end{tabular}
\end{center}

\noindent
\hl{Whenever $D_1$ detects the photon, $D_2$ does not detect that same photon, and vice versa.} This relationship is called \definition{Anticorrelation}. In other words, detecting the photon at one output \textbf{excludes} detecting the \textbf{same photon} at the other output.
\begin{gather*}
    D_1 \text{ clicks} \implies D_2 \text{ does not click}\\
    D_2 \text{ clicks} \implies D_1 \text{ does not click}
\end{gather*}
This is what we mean by ``\emph{anticorrelated detector clicks}''.

\highspace
\textcolor{Green3}{\faIcon{question-circle} \textbf{Why is this observation important?}} When a continuous light beam enters, the outgoing intensity is distributed equally:
\begin{equation*}
    I_1=\frac{I}{2},\qquad I_2=\frac{I}{2}
\end{equation*}
Both outputs carry light. Now consider a single photon. If the photon behaved like a \textbf{classical wave} whose energy was simply \textbf{divided into two equal portions}, we might expect half its energy to be detected at each output. But \hl{that's not what we observe!} Instead, the photon produces a localized detection at one output or the other.
\begin{equation*}
    P(D_1\text{ only}) = \frac{1}{2}, \qquad
    P(D_2\text{ only}) = \frac{1}{2}, \qquad
    P(D_1\text{ and }D_2) = 0
\end{equation*}
The third expression says that \hl{a single photon does not produce simultaneous detection events at both detectors.}

\highspace
\textcolor{Green3}{\faIcon{atom} \textbf{Particle indivisibility.}} Suppose a photon has energy:
\begin{equation*}
    E = h\nu
\end{equation*}
The two possible ideal detection outcomes are:
\begin{equation*}
    \left(E_1, E_2\right) =
    \begin{cases}
        (h\nu, 0) & \text{photon detected at } D_1 \\
        (0, h\nu) & \text{photon detected at } D_2
    \end{cases}
\end{equation*}
We do not observe the single photon's energy split into two separately detected halves. So \textbf{the photon is indivisible in detection}. But there is a subtle distinction worth preserving: anticorrelated \hl{clicks tell us \textbf{where the photon is detected}}, not necessarily which definite classical trajectory it followed \textbf{before} detection. Therefore, we must \textbf{not conclude} that the photon secretly chose one path before reaching the beam splitter. That is a possible classical-style hypothesis, not something established by anticorrelation alone.

\newpage

\noindent
\textcolor{Green3}{\faIcon{check-circle} \textbf{Anticorrelation does not contradict the 50/50 probabilities.}} Let's be precise about the difference between the two statements:
\begin{itemize}
    \item A \textbf{50/50 probability distribution} means that, after many repetitions, approximately half the photons are detected at $D_1$, $N_1\approx \frac{N}{2}$, and half at $D_2$, $N_2\approx\frac{N}{2}$.
    \item \textbf{Anticorrelation} means that for each individual photon, the two ideal detector outcomes are mutually exclusive (never both occur).
\end{itemize}
We can express this with conditional probabilities:
\begin{equation*}
    P(D_2\text{ clicks}\mid D_1\text{ clicks})=0,
    \qquad
    P(D_1\text{ clicks}\mid D_2\text{ clicks})=0
\end{equation*}
Considering the first conditional probability, this means that given that \hl{$D_1$ detected this photon, the probability that $D_2$ also detected the same photon is zero.} These two facts are completely compatible:
\begin{itemize}
    \item Over many photons, both detectors register approximately equal numbers of clicks (50/50 statistics).
    \item For each individual photon, exactly one detector clicks (anticorrelation).
\end{itemize}

\begin{figure}[!htp]
    \centering
    \tikzsetnextfilename{beam-splitter-anticorrelation-v3}
    \begin{tikzpicture}[
            font=\small,
            >={Stealth[length=2mm]},
            photon/.style={red!70!black},
            port/.style={photon, thick, dashed, ->},
            detector/.style={thick, draw=black!80, fill=gray!20},
            click/.style={draw=red!40!black, fill=red!70!black, minimum size=4.2mm, inner sep=0pt},
            silent/.style={draw=gray!60, fill=white, minimum size=4.2mm, inner sep=0pt},
            title/.style={anchor=west, font=\bfseries\small},
            note/.style={font=\scriptsize},
        ]
        % ================= (a) set-up and outcomes for ONE photon =================
        \node[title] at (0,2.75) {(a) One photon at a 50/50 beam splitter};

        % source
        \node[draw, thick, fill=yellow!20, align=center, font=\scriptsize,
        minimum height=8mm, inner xsep=2pt] (src) at (0.8,0) {single-photon\\source};

        % beam splitter: cube with a diagonal
        \coordinate (bs) at (3.0,0);
        \draw[thick, fill=cyan!10] ($(bs)+(-0.35,-0.35)$) rectangle ($(bs)+(0.35,0.35)$);
        \draw[very thick, cyan!50!black] ($(bs)+(-0.35,-0.35)$) -- ($(bs)+(0.35,0.35)$);
        \node[note, below=1mm] at ($(bs)+(0,-0.35)$) {50/50 BS};

        % input beam carrying one wave packet of energy h nu
        \draw[photon, thick] (src.east) -- ($(bs)+(-0.35,0)$);
        \draw[photon, thick, smooth, samples=80, domain=-0.4:0.4, shift={(2.25,0)}]
        plot (\x, {0.2*exp(-14*\x*\x)*sin(1900*\x)});
        \node[photon, above] at (2.25,0.22) {$h\nu$};

        % detectors (D-shaped, flat face towards the beam)
        \draw[detector] (4.85,-0.5) arc[start angle=-90, end angle=90, radius=0.5] -- cycle;
        \node[font=\footnotesize] at (5.09,0) {$D_1$};
        \draw[detector] (2.5,1.75) arc[start angle=180, end angle=0, radius=0.5] -- cycle;
        \node[font=\footnotesize] at (3.0,1.95) {$D_2$};

        % the two output ports: dashed = possible outcomes, not a known trajectory
        \draw[port] ($(bs)+(0.35,0)$) -- (4.83,0) node[midway, above, note, black] {output 1};
        \draw[port] ($(bs)+(0,0.35)$) -- (3.0,1.73) node[midway, right, note, black] {output 2};

        % ideal outcomes for each photon, drawn with the same click symbols as (b)
        \begin{scope}[xshift=6.3cm]
            \node[note, anchor=west] at (-0.2,2.25) {ideal outcomes per photon};
            \node[note] at (0.25,1.85) {$D_1$};
            \node[note] at (0.80,1.85) {$D_2$};
            % (h nu, 0)
            \node[click]  at (0.25,1.4) {};  \node[silent] at (0.80,1.4) {};
            \node[anchor=west] at (1.15,1.4) {$(h\nu,0)$};
            \node[anchor=west] at (3.75,1.4) {$P=\tfrac12$};
            % (0, h nu)
            \node[silent] at (0.25,0.75) {};  \node[click] at (0.80,0.75) {};
            \node[anchor=west] at (1.15,0.75) {$(0,h\nu)$};
            \node[anchor=west] at (3.75,0.75) {$P=\tfrac12$};
            % both click, or energy split in two: never observed
            \node[click, opacity=0.35] at (0.25,0.1) {};  \node[click, opacity=0.35] at (0.80,0.1) {};
            \draw[red!80!black, very thick] (0.0,-0.15) -- (1.05,0.35) (0.0,0.35) -- (1.05,-0.15);
            \node[anchor=west] at (1.15,0.1) {both, or $\tfrac{h\nu}{2}$ each};
            \node[anchor=west, red!80!black] at (3.75,0.1) {$P=0$};
            \node[note, anchor=west, align=left] at (-0.2,-0.45)
            {$P(D_1\text{ and }D_2)=0$: never both, never half};
        \end{scope}

        % ================= (b) click record for photons sent one at a time =================
        \begin{scope}[yshift=-2.8cm]
            \node[title] at (0,1.15) {(b) Click record, one photon at a time};

            \node at (0.45,0)     {$D_1$};
            \node at (0.45,-0.65) {$D_2$};
            \node[note, anchor=east] at (0.75,0.6) {photon};

            % illustrative sequence: 1 = D1 clicks, 2 = D2 clicks (photons 1 to 5 as in the notes)
            \foreach \w [count=\n] in {1,2,2,1,1,2,1,2,2,1} {
                \pgfmathsetmacro{\x}{0.45+0.6*\n}
                \node[note] at (\x,0.6) {\n};
                \ifnum\w=1
                \node[click]  at (\x,0) {};    \node[silent] at (\x,-0.65) {};
                \else
                \node[silent] at (\x,0) {};    \node[click]  at (\x,-0.65) {};
                \fi
            }

            % one column highlighted: anticorrelation
            \draw[blue!60!black, thick, rounded corners=2pt] (0.78,-0.97) rectangle (1.32,0.32);
            \draw[<-, blue!60!black] (1.05,-1.0) -- (1.05,-1.35)
            node[below, note, align=left, anchor=north west, xshift=-3mm, text=blue!60!black]
            {each column: exactly one click\\$P(D_2\text{ clicks}\mid D_1\text{ clicks})=0$};

            % row totals: 50/50 statistics
            \node[anchor=west] at (6.75,0)     {$N_1=5$};
            \node[anchor=west] at (6.75,-0.65) {$N_2=5$};
            \draw[decorate, decoration={brace, amplitude=4pt}] (7.95,0.2) -- (7.95,-0.85)
            node[midway, right=5pt, note, align=left]
            {$N_i\approx N/2$\\(50/50 statistics)};

            % legend
            \node[click]  at (6.95,-1.55) {};
            \node[note, anchor=west] at (7.15,-1.55) {click (whole $h\nu$)};
            \node[silent] at (9.8,-1.55) {};
            \node[note, anchor=west] at (10.0,-1.55) {no click};
        \end{scope}
    \end{tikzpicture}
    \caption{Anticorrelated detector clicks at a 50/50 beam splitter.
        (a)~A single photon of energy $h\nu$ is detected whole, at $D_1$ or at $D_2$,
        never at both and never split: $P(D_1\text{ and }D_2)=0$. Dashed lines mark the
        two possible output ports, not a trajectory the photon is known to follow.
        (b)~An illustrative sequence of detections. Read by columns, the record shows anticorrelation (one click per photon);
        read by rows, it shows the 50/50 statistics $N_1\approx N_2\approx N/2$
    (real runs fluctuate around $N/2$).}
    \label{fig:beam-splitter-anticorrelation-v3}
\end{figure}

\newpage

\begin{flushleft}
    \textcolor{Green3}{\faIcon{question-circle} \textbf{Why a beam splitter resembles a random 50/50 process}}
\end{flushleft}
We have established two important facts about a single photon entering an ideal 50/50 beam splitter:
\begin{enumerate}
    \item \textbf{Unpredictability}: we cannot predict which detector will register the photon.
    \item \textbf{Anticorrelation}: only one detector registers that photon, never both simultaneously.
\end{enumerate}
But now we have a natural question: \emph{\textbf{is a beam splitter simply tossing a coin for each photon?}}

\highspace
\textcolor{Green3}{\faIcon{coins} \textbf{The analogy with tossing a coin.}} Imagine an ordinary fair coin. Each toss produces one of two possible outcomes:
\begin{equation*}
    P(H)=\frac{1}{2},\qquad P(T)=\frac{1}{2}
\end{equation*}
Similarly, a photon entering a 50/50 beam splitter produces one of two possible detection outcomes:
\begin{equation*}
    P(D_1)=\frac{1}{2},\qquad P(D_2)=\frac{1}{2}
\end{equation*}
If we repeat either experiment $1'000$ times, we expect approximately $500$ occurrences of each outcome. From the perspective of \textbf{individual results and their frequencies}, the \hl{two experiments look remarkably similar.} But does that mean they are governed by the same kind of randomness?

\highspace
\textcolor{Green3}{\faIcon{question-circle} \textbf{Could the photon behave like a classical coin?}} In classical mechanics, a coin toss appears random because we do not know the exact initial conditions (remark \autopageref{def:epistemic-randomness}). \emph{What if the photon behaves similarly?} Perhaps the photon has some unknown microscopic properties that determine whether it will ultimately be detected at $D_1$ or $D_2$.

\highspace
For example, we could hypothesize:
\begin{equation*}
    \begin{gathered}
        \text{Photon enters the beam splitter}\\
        \Downarrow\\
        \text{Some unknown physical details determine the outcome}\\
        \Downarrow\\
        D_1\text{ or }D_2
    \end{gathered}
\end{equation*}
In this hypothetical picture, the 50/50 statistics arise because we lack information about those microscopic details. This is the \definition{Classical-Ignorance Hypothesis}: \hl{an event appears random only because we do not know all the physical information that determines its outcome.}

\highspace
\textcolor{Green3}{\faIcon{question-circle} \textbf{Can the single beam-splitter experiment distinguish the two explanations?}} Similar to the single-particle regime of the double-slit experiment, \hl{a single 50/50 beam splitter cannot, by itself, distinguish quantum probabilities from a classical coin-toss model.} All we have established is that its outcomes \textbf{behave statistically like a fair coin toss}. Whether the underlying physics is the same remains an open question.

\highspace
\textcolor{Green3}{\faIcon{question-circle} \textbf{How can we investigate the difference?}} We need an experiment in which the classical random-choice model and quantum mechanics make \textbf{different predictions}.

\highspace
Instead of studying one beam splitter in isolation, what if we connect two beam splitters using mirrors, allowing the two possible routes to meet again? This leads us to the \textbf{Mach-Zehnder interferometer}. The idea is really simple: \hl{if each beam splitter acts like an independent coin toss, combining two of them should still produce the output probabilities predicted by those random choices.} In other words, if each beam splitter behaves like a fair coin toss, then passing a photon through two beam splitters should be equivalent to tossing a coin twice.

\highspace

\begin{takeawaysbox}
    \textbf{A single photon at a 50/50 beam splitter.}

    An ideal 50/50 beam splitter produces two equally
    probable detection outcomes:
    \begin{equation*}
        P(D_1)=P(D_2)=\frac{1}{2}
    \end{equation*}
    However, \hl{one photon produces one complete
    detection event}, either at $D_1$ or at $D_2$,
    never at both simultaneously (anticorrelation).
    The photon is \textbf{indivisible in detection}.

    \highspace
    \textbf{The classical coin-toss analogy.}

    The 50/50 statistics resemble tossing a fair coin.
    Perhaps the outcome is determined by unknown
    microscopic information, as in classical mechanics
    (the \emph{classical-ignorance hypothesis}).

    \highspace
    However, \hl{identical probabilities do not imply
    identical physical mechanisms}. A single beam
    splitter cannot distinguish between an ordinary
    classical random-choice model and the quantum
    description.

    \highspace
    \textbf{The Mach-Zehnder interferometer.}
    If a beam splitter really
    behaves like a classical coin toss, should two
    successive beam splitters behave like two independent
    coin tosses? The Mach-Zehnder interferometer will test this idea.
\end{takeawaysbox}
